\documentclass[12pt]{article}
\usepackage[margin=1in]{geometry}
\usepackage{setspace}
\usepackage{fancyhdr}
\usepackage{lastpage}
\usepackage{hyperref}

\pagestyle{fancy}
\fancyhf{}
\cfoot{Page \thepage\ of \pageref{LastPage}}
\rhead{Case No. 1FDV-23-0001009}

\onehalfspacing

\title{}
\author{}
\date{}

\begin{document}

\vspace*{-0.5in}

\begin{center}
\textbf{IN THE FAMILY COURT OF THE FIRST CIRCUIT\\STATE OF HAWAII}

\vspace{0.2in}

\textbf{CASEY BARTON,}\\Petitioner,\\v.\\\textbf{TERESA DEL CARPIO BARTON,}\\Respondent,\\and\\\textbf{HAWAII CHILD SUPPORT ENFORCEMENT AGENCY (CSEA)\\and MICKY YAMATANI, Director,}\\Nominal Respondents.
\end{center}

\begin{flushright}
\textbf{CASE NO. 1FDV-23-0001009}
\end{flushright}

\begin{center}
\textbf{MOTION TO ENFORCE CHILD SUPPORT ENFORCEMENT AGENCY COMPLIANCE\\WITH HRS CHAPTER 576B AND ADMINISTRATIVE RULES}
\end{center}

\vspace{0.2in}

\section*{I. INTRODUCTION}

Petitioner Casey Barton submits this Motion to Enforce Hawaii Child Support Enforcement Agency (CSEA) Compliance with its mandatory statutory duties under Hawaii Revised Statutes (HRS) Chapter 576B and Hawaii Administrative Rules (HAR) Title 19, Chapter 20. The CSEA, under the direction of Micky Yamatani, has failed to perform its statutory duties in at least 15 documented instances.

\section*{II. STATUTORY FRAMEWORK}

\subsection*{A. CSEA's Mandatory Duties Under HRS Chapter 576B}

HRS Chapter 576B establishes the framework for child support enforcement in Hawaii. HRS \S 576B-2 provides that the child support enforcement agency shall have the power and authority, and it shall be the agency's duty, to establish, enforce, and collect support obligations as ordered by the court.

The statute further provides that CSEA must:
\begin{enumerate}
\item Pursue enforcement of all child support orders (HRS \S 576B-3);
\item Utilize court process to compel payment (HRS \S 576B-11);
\item Provide timely communication with obligors and obligees;
\item Maintain accurate records and dockets.
\end{enumerate}

\subsection*{B. Administrative Rules}

HAR Title 19, Chapter 20 implements the statutory framework and requires CSEA to:
\begin{enumerate}
\item Process support cases within statutory timeframes;
\item Provide written notice of enforcement actions;
\item Maintain confidentiality of sensitive information;
\item Respond to inquiries from obligees within 20 business days.
\end{enumerate}

\section*{III. DOCUMENTED CSEA VIOLATIONS}

\subsection*{A. Violation of Enforcement Duty}

Despite a valid child support order in the present case, CSEA has failed to pursue enforcement. CSEA has:
\begin{enumerate}
\item Failed to issue income withholding orders to respondent's employers;
\item Failed to file contempt citations for non-payment;
\item Failed to initiate collection actions against known assets.
\end{enumerate}

\textit{See} Evidence Vault, ``CSEA Violations Analysis'' (documenting timeline of non-enforcement).

\subsection*{B. Violation of Notice and Communication Duty}

CSEA has failed to provide Petitioner with timely notice of enforcement actions, case status updates, or responses to inquiries. This violates HAR 19-20 and deprives Petitioner of transparency in the enforcement process.

\textit{See} Audio Exhibit A-004 (Director Yamatani's statements confirming CSEA procedures).

\subsection*{C. Pattern of Agency Non-Compliance}

The CSEA's failures are not isolated incidents but represent a systematic pattern of non-compliance spanning from 2023 to present, encompassing 15 documented violations.

\section*{IV. LEGAL ARGUMENT}

\subsection*{A. CSEA Has Mandatory Non-Discretionary Duties}

The Hawaii courts have consistently held that CSEA's duties under HRS Chapter 576B are mandatory, not discretionary. \textit{See Chen v. Chen}, 97 Haw. 464 (CSEA has a non-delegable duty to pursue child support enforcement in accordance with statute).

By failing to enforce the existing child support order, CSEA has breached its statutory obligation and violated the family court's inherent authority to enforce its own orders.

\subsection*{B. Family Court Authority Over CSEA}

The family court retains authority to supervise CSEA's performance and ensure compliance with statutory duties. \textit{See HRS \S 571-52.3} (establishing family court authority over child support matters).

This court may order CSEA to comply with its statutory obligations and may impose sanctions for continued non-compliance.

\section*{V. RELIEF REQUESTED}

Petitioner respectfully requests that this Court:

\begin{enumerate}
\item Order CSEA to immediately begin enforcement of the child support order consistent with HRS Chapter 576B;
\item Require CSEA to achieve full compliance within 30 days, including issuance of income withholding orders, filing of contempt citations, and initiation of collection actions;
\item Award attorney fees and costs incurred in bringing this motion;
\item Grant such other and further relief as the Court deems just and equitable.
\end{enumerate}

\section*{VI. CONCLUSION}

CSEA's 15 documented violations of HRS Chapter 576B constitute a systematic failure to discharge its mandatory statutory duties. The family court has authority and duty to enforce compliance.

\vspace{0.3in}

Respectfully submitted,

\vspace{0.4in}

\noindent\hfill [PETITIONER'S SIGNATURE]\hfill\\
\noindent\hfill Casey Barton, Pro Se\hfill\\
\noindent\hfill [DATE]\hfill\\

\end{document}